\section{Results}
\label{sec:results}

\subsection{Validation of the SAR area and water-level inputs}
\label{sec:res_inputvalid}
The Sentinel-1 water-surface-area input carries the reliability \citet{PaperOne} establishes
across 66 reservoirs, including 5 Sicilian ones. Surface-area skill correlates with shoreline
area-to-perimeter ratio (Spearman $\rho=0.51$), reaching a median KGE of $0.93$ at high A/P
($\geq200$~m) against $0.48$ at low A/P ($<100$~m), with $82\%$ of that global set clearing
KGE $\geq0.5$. We take that reliability as given rather than re-validating it here.

\label{sec:res_ladder}
For the water level, Poma offers the clearest single case, showing the textbook signature of
a pure offset rather than a real water-level disagreement, with a correlation of $r=1.00$
before correction and the mean offset alone accounting for essentially all of the raw RMSE.
Correcting it changes RMSE, which is sensitive to an absolute offset, but leaves KGE (whose
mean-ratio term is small relative to these reservoirs' absolute elevations) essentially
unaffected. Fig.~\ref{fig:wlsource} plots the cleaned SWOT against the gauge over the full
SWOT era (2023-07 onward) for all nine reservoirs, and Table~\ref{tab:wlrmse} summarises the
agreement as both an RMSE and a KGE (Section~\ref{sec:inputvalid}). We call a source
\emph{acceptable} when its error is small relative to the reconstruction's own elevation bias
(a few metres in amplitude, Section~\ref{sec:res_garcia}), requiring KGE $\geq0.8$ once cleaned.
Seven of nine reservoirs clear this bar: Arancio, Poma, Garcia, Rosamarina, Olivo, Castello
and Ancipa, RMSE $0.14$--$0.63$~m, KGE $0.94$--$0.99$.

\begin{table}[htbp]\centering\small
\caption{SWOT versus in-situ gauge water level, cleaned and with malfunction windows
excluded, before any volume is computed, ordered by A/P (Fig.~\ref{fig:wlsource} shows the
underlying series). Shift is the constant offset removed from SWOT, and $r$, $\alpha$,
$\beta$ are the KGE components (Section~\ref{sec:inputvalid}). Tier ``truth'' means an
independent survey or curve exists for the resulting band volume
(Section~\ref{sec:res_curves}); ``consistency'' means gauge and SWOT are compared only to
each other.}
\label{tab:wlrmse}
\begin{tabular}{lrrrrrrrl}
\toprule
Reservoir & A/P (m) & WL RMSE (m) & Shift (m) & $r$ & $\alpha$ & $\beta$ & KGE & Tier \\
\midrule
Olivo       &  50.7 & 0.26 & $-0.74$ & 0.99 & 1.00 & 1.00 & 0.99 & truth \\
Ancipa      &  90.5 & 0.38 & $-2.72$ & 1.00 & 1.04 & 1.00 & 0.96 & consistency$^c$ \\
Nicoletti   & 119.7 & 0.76 & $+1.89$ & 0.98 & 1.32 & 1.00 & 0.68 & truth$^b$ \\
Castello    & 126.7 & 0.63 & $+0.04$ & 0.99 & 1.04 & 1.00 & 0.96 & truth \\
Garcia      & 167.7 & 0.14 & $-1.55$ & 1.00 & 0.99 & 1.00 & 0.99 & truth \\
Arancio     & 182.2 & 0.59 & $+0.21$ & 0.99 & 0.94 & 1.00 & 0.94 & truth \\
Rosamarina  & 187.4 & 0.60 & $+0.28$ & 1.00 & 0.98 & 1.00 & 0.98 & truth \\
Poma        & 190.1 & 0.16 & $+1.90$ & 1.00 & 0.99 & 1.00 & 0.99 & truth \\
Pozzillo    & 240.5 & 1.77 & $-1.56$ & 1.00 & 1.62 & 1.00 & 0.38 & consistency$^a$ \\
\bottomrule
\end{tabular}

\vspace{2pt}
{\footnotesize $^a$Pozzillo's low KGE is a variance mismatch ($\alpha=1.62$), not a real
disagreement ($r=1.00$). $^b$Nicoletti's tier holds once its gauge freeze is excluded, with a
residual caveat. $^c$Ancipa's agreement is conditional on excluding three dates. All three are
discussed in the text below.}
\end{table}

\begin{figure}[htbp]\centering
  \includegraphics[width=\linewidth]{wl_source_grid.png}
  \caption{Gauge versus SWOT water level, SWOT era only (2023-07 onward), all nine
  reservoirs, ordered by A/P. Grey markers are SWOT points removed by outlier cleaning;
  shaded bands are gauge malfunction windows (excluded, shown faded). RMSE and KGE
  annotated for all nine.}
  \label{fig:wlsource}
\end{figure}

At the low-A/P Ancipa, SWOT at first appeared to diverge sharply from the gauge. This traced
to three specific dates (2024-10-18, 2024-11-08, 2024-11-29) where SWOT read a nearly-constant
${\sim}920$--$924$~m while the gauge, and independently the SAR-measured water area, both
showed a real, continuous drawdown toward the reservoir's minimum (SAR area falling to
$0.26$~ha by mid-November, essentially dry). SWOT's Prior Lake detection did not track the
water body once it shrank to a sliver, and the raw
SWOT record has no observations at all for the following five months (2024-11-29 to
2025-04-24), a genuine gap in the product itself, not one introduced by our cleaning.
Excluding these three dates, gauge and SWOT agree almost perfectly ($r=1.00$,
Table~\ref{tab:wlrmse}), so SWOT is reliable at Ancipa away from its most extreme drawdown.
Olivo,
Castello and Nicoletti also have usable SWOT records, once matched to each reservoir's own
footprint rather than its granule's metadata, reported above. DAHITI is not evaluated
separately. For the reservoirs where it exists it is a reprocessed
SWOT product, matching SWOT epochs at $>96\%$ of dates to within $\pm0.6$~m.

Four of the nine gauges have a documented reliability problem in the publicly available
record we rely on (Section~\ref{sec:data}), a limitation of what is accessible to us rather
than necessarily of the water authority's own internal instrumentation; each is resolved with
the same substitution logic rather than by discarding the affected dates: Rosamarina's sensor
froze at two different flat readings across
$\sim$19 months combined; Garcia's sensor reads the dry lakebed floor instead of the true level
whenever the reservoir empties below $\sim$176~m, a fault that recurred twice a year apart
(2024-09-04 to 2025-02-13 and 2025-08-06 to 2026-02-03); Nicoletti's sensor was inoperative
until 2023-11 and then, once it began reporting, frozen at a single flat value for a further 14
months (2023-11-15 to 2025-01-17). Nicoletti's gauge-SWOT offset also shifts across the record,
from a tight $-1.89$~m (2025-02 to 2026-02) to $-0.2$ to $-0.6$~m from 2026-03 on; the cause is
unexplained, and we leave it as a disclosed limitation rather than a second correction. Pozzillo's
sensor combines both patterns, with a 13-month frozen episode (2023-08-27 to 2024-10-03) on top
of a persistent range-compression across the rest of the record, caught only because an
independent signal (SAR water area) was checked against both candidate water levels rather than
the gauge being trusted by default. Gauge and SWOT correlate almost perfectly ($r=1.00$) once
Pozzillo's frozen episode is excluded; the remaining KGE gap is a pure variance mismatch
($\alpha=1.62$), and the SAR area corroborates SWOT independently, correlating more strongly
with SWOT ($r=0.90$) than with the gauge ($r=0.76$) over the same period. The gauge under-represents
Pozzillo's true range; this is not a SWOT error. Rosamarina's, Garcia's and Nicoletti's are
intermittent malfunctions with a clear onset and recovery; Pozzillo's range-compression is a
persistent, whole-record limitation layered on top of its own frozen episode. In
each case SWOT (Rosamarina, Garcia, Pozzillo) fills the gap for the affected reconstruction
dates; Nicoletti's own selected reconstruction dates happen to all postdate its gauge's
recovery, so its freeze affects the water-level validation above but not its reconstruction,
which the boletin-derived fallback already covers for dates before the gauge existed at all
(Section~\ref{sec:data}). That four of nine gauges needed this kind of substitution in
practice, not merely in principle, is itself evidence for pairing SAR with a gauge-independent
altimetry source rather than relying on in-situ networks alone. What this water-level
substitution does to reconstructed \emph{volume}, once it is carried through the
reconstruction, is a question about the reconstruction's output, not the input, and is
addressed next (Section~\ref{sec:res_hyps}).

\subsection{Reconstructed bathymetry and capacity change}
\label{sec:res_change}
With both inputs validated, the reconstruction follows in three steps.
Section~\ref{sec:res_hyps} compares the hypsometric curve implied by each water-level
source, including one built with no gauge anywhere in the chain.
Section~\ref{sec:res_dem} reports the bathymetric DEM and capacity-change estimate that
follow from those curves, and how much removing the gauge alters the estimate.
Section~\ref{sec:res_voltime} then applies the reconstructed curve to each reservoir's
full historical record and scores it against the water authority's own registered volumes.

\subsubsection{Hypsometric curve: gauge+SWOT-fallback versus full remote sensing}
\label{sec:res_hyps}
For each reservoir we fit the same power-law hypsometric curve $A=a(h-h_0)^b$ two ways: (a) on
the gauge+SWOT calibration pairs (gauge, substituting SWOT only inside a documented
malfunction window, Section~\ref{sec:res_inputvalid}); and (b) on genuine SAR-area/SWOT-level
coincident pairs alone, with no gauge value anywhere in the fit, which gives a
full-remote-sensing (FRS) reconstruction. Both are compared, over the water-level range they share, against (c)
the best available independent reference (an updated survey or curve, or the original design
curve where none exists). Fig.~\ref{fig:hyps} plots all three curves against the real
calibration pairs for all nine reservoirs; Table~\ref{tab:hyps} reports the resulting band
volumes.

\begin{figure}[htbp]\centering
  \includegraphics[width=\linewidth]{hypsometry_comparison.png}
  \caption{Hypsometric curve by water-level source: (a) gauge+SWOT-fallback,
  (b) full remote sensing (SWOT-only, no gauge), and (c) the best available independent
  reference, for all nine reservoirs, ordered by A/P. Grey points are the gauge+SWOT
  calibration pairs.}
  \label{fig:hyps}
\end{figure}

\begin{table}[htbp]\centering\small
\caption{Band volume (Mm$^3$) by hypsometric curve. The comparable band is the
water-level range over which curves (a) and (b) can both be evaluated, that is the
overlap of the two sets of calibration pairs. It is narrower than the reconstructed band
of Table~\ref{tab:capacity} wherever one source covers less of the drawdown than the
other, and is not the range the DEM spans. Reference = updated survey/curve, except
Ancipa and Pozzillo (design curve; no updated survey exists for either). $n_b$ = number
of genuine SAR-area/SWOT-level pairs the FRS fit draws on (unlike (a), not fixed at 10).
$^{\dagger}$Ancipa is the extreme case on both counts. Its raw SWOT record is sparse (24
observations, against 83--145 elsewhere) and SWOT stops tracking the reservoir below
${\sim}920$~m (Section~\ref{sec:res_inputvalid}), so its comparable band covers only the
upper $25$~m even though the reconstruction itself spans $898.5$--$945.1$~m. Pairing is
not at fault: $67\%$ of its quality-screened SWOT passes find a SAR area within $\pm3$
days, in line with the other reservoirs.}
\label{tab:hyps}
\resizebox{\linewidth}{!}{%
\begin{tabular}{lrrrrrrr}
\toprule
Reservoir & A/P (m) & Comparable band (m) & $n_b$ & Vol.\ (a) & Vol.\ (b) & Vol.\ (c) & (b$-$a) \\
\midrule
Olivo       &  50.7 & 432.1--439.0 & 98 &  2.09 &  2.68 &  3.45 & $+0.59$ \\
Ancipa      &  90.5 & 920.3--945.2 & 12\rlap{$^{\dagger}$} & 12.03 & 12.82 & 19.38 & $+0.79$ \\
Nicoletti   & 119.7 & 375.1--379.9 & 50 &  3.22 &  3.55 &  5.02 & $+0.33$ \\
Castello    & 126.7 & 277.5--292.8 & 49 & 13.66 & 13.63 & 14.90 & $-0.03$ \\
Garcia      & 167.7 & 171.0--189.7 & 53 & 41.78 & 42.41 & 41.09 & $+0.63$ \\
Arancio     & 182.2 & 167.2--176.4 & 69 & 15.23 & 15.60 & 17.15 & $+0.37$ \\
Rosamarina  & 187.4 & 140.6--162.1 & 45 & 45.92 & 46.65 & 49.31 & $+0.72$ \\
Poma        & 190.1 & 176.5--193.2 & 51 & 41.80 & 42.05 & 46.87 & $+0.25$ \\
Pozzillo    & 240.5 & 342.9--354.9 & 25 & 44.57 & 45.91 & 46.49 & $+1.34$ \\
\bottomrule
\end{tabular}}
\end{table}

The FRS fit draws on far more genuine coincident pairs than the fixed 10-date gauge+SWOT
calibration set ($n_b=12$--$98$), since it is not tied to the specific dates exported for the
reconstruction. Its band volume differs from the gauge+SWOT curve's by only
$0.03$--$1.34$~\si{Mm^3} at every reservoir (Table~\ref{tab:hyps}), consistently smaller
than either curve's own gap to the independent reference, which reaches $0.58$--$7.35$~\si{Mm^3}
and does not track A/P as cleanly here as the capacity-change comparison in
Section~\ref{sec:res_dem} does. Garcia is the one reservoir where both curves sit
\emph{above} its echo-sounder-referenced curve ($+0.69$/$+1.32$~\si{Mm^3}) rather than below
it, though by a margin no larger than most other reservoirs' own gaps. Removing the gauge
entirely, in other words, changes the reconstructed hypsometry by less than the
reconstruction's own disagreement with independent ground truth, at all nine reservoirs.

\subsubsection{Bathymetric reconstruction and capacity change}
\label{sec:res_dem}
The waterline-stacked DEMs reconstruct the exposed basin of all nine reservoirs over the
2022--2026 window, yielding 2D/3D bathymetry and a capacity-change estimate against each
reservoir's design curve, for both the gauge+SWOT-fallback and the FRS water-level
sources (Table~\ref{tab:capacity}, Fig.~\ref{fig:change}). Pozzillo's gauge+SWOT DEM is built
from SWOT rather than the gauge, per Section~\ref{sec:res_inputvalid}; every other reservoir's
gauge+SWOT DEM uses the gauge (substituting SWOT or the design curve only inside a documented
malfunction window). Relative to the design curve, the gauge+SWOT SAR-observable (band)
capacity change ranges from a modest $-0.4\%$ (Pozzillo) to $-45.3\%$ (Olivo). The exposed
band itself represents a widely varying share of each reservoir's total design-curve volume,
$39.3\%$--$92.0\%$ (mean $70.6\%$ across the eight reservoirs for which this ratio is well
posed), so for most reservoirs the majority of storage is directly observed, but the
fraction is far from uniform and is smallest exactly where A/P is lowest. Ancipa is excluded
from this ratio, because its design curve's lowest tabulated quota sits above the DEM's true
floor, so the raw ratio exceeds 100\% (a design-curve extrapolation artifact, not a real deep
zone); we report it as $100\%$ with this caveat rather than force a spurious number.

\begin{table}[htbp]\centering\small
\caption{Reconstructed capacity change by reservoir, ordered by A/P, for the two
water-level variants (a) gauge+SWOT-fallback and (b) FRS (SWOT-only). Band-observable is
the share of design-curve volume the exposed band represents. SAR band change is measured
relative to the DEM floor. The independent columns give the updated survey or curve for
variant (a) only (Section~\ref{sec:res_curves}), with \emph{total} including the unobserved
deep zone. $^a$Capped at 100\% (design-curve extrapolation artifact, see text).
$^b$Pozzillo's variant (a) uses SWOT as its primary water level.}
\label{tab:capacity}
\resizebox{\linewidth}{!}{%
\begin{tabular}{lrrrrrrr}
\toprule
& & \multicolumn{2}{c}{Band-observable (\%)} & \multicolumn{2}{c}{SAR band change (\%)} & \multicolumn{2}{c}{Independent (\%)} \\
\cmidrule(lr){3-4}\cmidrule(lr){5-6}\cmidrule(lr){7-8}
Reservoir & A/P (m) & (a) & (b) & (a) & (b) & band & total \\
\midrule
Olivo        &  50.7 & 50.1 & 51.0 & $-45.3$ & $-29.6$ & $-3.6$  & $-31.2$ \\
Ancipa       &  90.5 & 100$^a$ & 100$^a$ & $-39.6$ & $-33.6$ & --- & --- \\
Nicoletti    & 119.7 & 39.3 & 49.6 & $-40.2$ & $-34.9$ & $-6.1$  & $-5.0$ \\
Castello     & 126.7 & 67.2 & 70.9 & $-30.0$ & $-29.8$ & $-23.5$ & $-31.9$ \\
Garcia       & 167.7 & 88.7 & 90.8 & $-20.8$ & $-20.2$ & $-22.1$ & $-30.9$ \\
Arancio      & 182.2 & 72.2 & 76.1 & $-15.5$ & $-14.4$ & $-3.7$  & $-15.5$ \\
Rosamarina   & 187.4 & 76.1 & 77.6 & $-15.1$ & $-12.3$ & $-7.3$  & $-16.7$ \\
Poma         & 190.1 & 78.8 & 79.2 & $-12.8$ & $-11.3$ & $-0.5$  & $+2.9$ \\
Pozzillo$^b$ & 240.5 & 92.0 & 93.4 & $-0.4$  & $-2.5$  & --- & --- \\
\bottomrule
\end{tabular}}
\end{table}

\begin{figure}[htbp]\centering
  \includegraphics[width=\linewidth]{bathymetry_change_summary.png}
  \caption{SAR-detected (gauge+SWOT) capacity change versus the design curve for all nine
  reservoirs (band-relative), with independent survey/curve change and, where available, the
  total loss including the deep zone. Ordered by A/P.}
  \label{fig:change}
\end{figure}

Beyond the capacity-change estimate, the spatial reconstruction itself is comparably robust
to removing the gauge. Rebuilding each reservoir's DEM entirely from SWOT reproduces the
gauge+SWOT DEM to below $1$~m RMSE at six of the nine (Arancio $0.29$~m, Poma $0.31$~m, Garcia
$0.50$~m, Castello $0.60$~m, Pozzillo $0.62$~m, Rosamarina $0.71$~m) and to $2.0$--$3.8$~m at
the remaining three (Ancipa $1.97$~m, Nicoletti $2.09$~m, Olivo $3.82$~m), precisely the
three lowest-A/P reservoirs in the set, with Olivo the clear outlier. Ancipa's agreement is
nonetheless better than Nicoletti's despite its lower A/P: its reconstruction combines two
Sentinel-1 orbits rather than the single orbit used at every other reservoir, which resolves a
near-dam radar-shadow gap in one otherwise-anomalous calibration date and, as a side effect,
roughly doubles its usable mask density.

The apparent band loss itself tracks A/P: the three lowest-A/P reservoirs (Olivo $50.7$~m,
Ancipa $90.5$~m, Nicoletti $119.7$~m) show the largest losses ($-45.3\%$, $-39.6\%$,
$-40.2\%$ respectively), while the two highest-A/P reservoirs (Poma $190.1$~m, Pozzillo
$240.5$~m) show the smallest ($-12.8\%$, $-0.4\%$). Fig.~\ref{fig:aevgrid} shows why directly:
at low A/P (Olivo, Nicoletti) the reconstructed volume--elevation curve sits well to the left
of \emph{both} the design and the updated curve at every level, not only at the top, while
the two reference curves stay close to each other. At medium/high A/P
(Garcia, Arancio, Rosamarina, Poma, Pozzillo) the SAR curve instead tracks the updated
curve closely across most of the range. Section~\ref{sec:res_curves} quantifies this pattern
directly.

\begin{figure}[htbp]\centering
  \includegraphics[width=\linewidth]{aev_grid.png}
  \caption{Reconstructed volume--elevation curve (gauge+SWOT, SAR) against the design curve
  and, where available, the updated official curve or echo-sounder survey, for all nine
  reservoirs, ordered by A/P.}
  \label{fig:aevgrid}
\end{figure}

\subsubsection{Volume time series against official monitoring records}
\label{sec:res_voltime}
Beyond validating the reconstruction itself, we test whether it improves day-to-day reservoir
management, applying two curves, the FRS hypsometric curve (Section~\ref{sec:res_hyps}) and
the best available independent reference (updated survey or curve, else design,
Section~\ref{sec:res_curves}), to each reservoir's full, multi-year Sentinel-1 area time
series, and comparing the resulting volume time series against the water authority's own
officially registered volume record (monthly, 2007--2026; daily, 2023--2026 where available),
independent of any curve used in this study. This covers all nine reservoirs. Five have a full
Sentinel-1 area record from the original export (Ancipa, Poma, Pozzillo, Rosamarina, Garcia),
and the remaining four (Arancio, Castello, Nicoletti, Olivo) were re-exported over the same
full Sentinel-1 era for this comparison.

Each curve only covers the area range it was itself fitted or tabulated over, so applying it
outside that range is extrapolation, not a genuine estimate; comparisons are restricted to
dates within each curve's own range, tracked separately since the FRS curve's own coverage
($39$--$98\%$ of each reservoir's record) can differ from the reference curve's
(Table~\ref{tab:opvol}). All nine area series enter this comparison as the same product,
carrying the outlier rejection and LOWESS smoothing of \citet{PaperOne}'s published
post-processing, which suppresses per-scene classification noise, not real volume
change, and scoring a smoothed series against an unsmoothed one would otherwise favour the
former. The hypsometric fits of Section~\ref{sec:res_hyps} instead use the raw
per-acquisition areas, since they pair individual SAR acquisitions with individual SWOT
passes, and smoothing would blend neighbouring dates into pairs never actually observed.

\begin{table}[htbp]\centering\small
\caption{FRS (SWOT-only) versus the best available independent reference curve, converting
each reservoir's own SAR area record to volume and comparing to the officially registered
volume, restricted to dates within each curve's own observed area range (coverage). Skill is
scored with KGE (Section~\ref{sec:inputvalid}), with RMSE alongside it. Unlike the
water-level comparison of Table~\ref{tab:wlrmse}, these are volumes rather than absolute
elevations, so the mean-ratio component is informative here.}
\label{tab:opvol}
\resizebox{\linewidth}{!}{%
\begin{tabular}{llrrrrr}
\toprule
& & & \multicolumn{2}{c}{FRS (SWOT-only)} & \multicolumn{2}{c}{Reference curve} \\
\cmidrule(lr){4-5}\cmidrule(lr){6-7}
Reservoir & Record & Coverage & RMSE (Mm$^3$) & KGE & RMSE (Mm$^3$) & KGE \\
\midrule
Olivo      & monthly & 53\% & 0.63 & \textbf{0.47} & 2.05  & 0.17 \\
Ancipa     & monthly & 39\% & 1.93 & \textbf{0.74} & 7.26  & 0.27 \\
Ancipa     & daily   & 39\% & ---  & ---           & 10.30 & $-0.01$ \\
Nicoletti  & monthly & 42\% & 2.46 & 0.49          & 2.12  & \textbf{0.64} \\
Castello   & monthly & 98\% & 3.95 & 0.70          & 3.04  & \textbf{0.72} \\
Castello   & daily   & 98\% & 4.00 & 0.43          & 2.26  & \textbf{0.65} \\
Garcia     & monthly & 94\% & 9.05 & \textbf{0.77} & 18.36 & 0.52 \\
Garcia     & daily   & 94\% & 3.94 & \textbf{0.68} & 9.16  & 0.28 \\
Arancio    & monthly & 78\% & 1.95 & \textbf{0.73} & 5.84  & 0.59 \\
Rosamarina & monthly & 81\% & 5.71 & \textbf{0.89} & 9.18  & 0.78 \\
Rosamarina & daily   & 81\% & 2.53 & \textbf{0.71} & 7.02  & 0.53 \\
Poma       & monthly & 95\% & 6.88 & \textbf{0.78} & 9.06  & 0.75 \\
Poma       & daily   & 95\% & 5.87 & \textbf{0.47} & 6.60  & 0.30 \\
Pozzillo   & monthly & 73\% & 5.97 & 0.80          & 6.95  & \textbf{0.87} \\
\bottomrule
\end{tabular}}
\end{table}

The FRS curve tracks the officially registered volume better than the independent reference
at six of the nine reservoirs, by KGE: Rosamarina ($0.89$ vs.\ $0.78$), Poma ($0.78$ vs.\
$0.75$), Garcia ($0.77$ vs.\ $0.52$), Ancipa ($0.74$ vs.\ $0.27$), Arancio ($0.73$ vs.\ $0.59$)
and Olivo ($0.47$ vs.\ $0.17$). The ranking is largely insensitive to the metric, since RMSE
favours the FRS curve at the same six reservoirs and at Pozzillo besides. What the KGE
decomposition adds is the reason. The reference curves fail mainly through their mean ratio,
which runs from $0.47$ to $0.86$ at every reservoir except Pozzillo, meaning they
systematically place storage below the registered record, as an outdated curve would. The FRS
curve's own mean ratio sits near unity ($0.83$--$1.16$) everywhere except Nicoletti ($1.39$)
and Castello ($1.29$).

The reference curve is not beaten at Nicoletti (KGE $0.49$ against $0.64$), at Castello
($0.70$ against $0.72$, effectively tied) or at Pozzillo ($0.80$ against $0.87$). Nicoletti and
Castello are also the two reservoirs whose FRS curve is calibrated over the narrowest slice of
their observed water-level range, the upper $29\%$ and $70\%$ respectively against $85\%$ at
Arancio, and both carry the largest FRS mean ratios in the set ($1.39$ and $1.29$). At Castello
the registered volume falls between the two curves, with FRS above it and the reference below.
At Pozzillo the two are close on RMSE ($5.97$ against $6.95$) and separate on the mean ratio,
$1.16$ for the FRS curve against $1.00$ for the reference.

\begin{figure}[htbp]\centering
  \includegraphics[width=\linewidth]{area_volume_timeseries.png}
  \caption{Volume time series from the Sentinel-1 area record, converted via the FRS (solid
  purple) and independent-reference (dashed green) curves, against the officially registered
  volume (points), for all nine reservoirs, ordered by A/P. Each curve is drawn solid
  only within its own observed area range; the faint dotted extension beyond it is
  extrapolation, excluded from Table~\ref{tab:opvol}.}
  \label{fig:opvol}
\end{figure}

\subsection{Validation against independent references}
\label{sec:res_valid}

\subsubsection{Field validation against an echo-sounder survey (Garcia)}
\label{sec:res_garcia}
The December-2025 echo-sounder survey supports two comparisons of different kinds.
Pixel-wise elevation error is measured against the survey gridded from its own raw
sonar and shore point cloud, since a per-pixel comparison requires a raster. All
volume comparisons instead use the curated official quota-area-volume table derived
from the same survey, because that raw grid interpolates sparse shore points and so
under-represents area with height above the survey's own waterline. Against it, the
Garcia DEM has an overall bias of $-0.88$~m and RMSE $2.67$~m in the SAR-observable
range ($\mathrm{R}^2=0.61$). A residual
elevation dependence remains, from $-1.58$~m near the floor ($175$--$176$~m) to
$-2.39$~m near the top ($189$--$190$~m) (Fig.~\ref{fig:garcia}). Band-relative to
design, the curated curve shows a true loss of $-22.1\%$ (Table~\ref{tab:capacity}),
of which the SAR DEM detects $-20.8\%$, capturing $93.7\%$ of the true band loss.
The remaining $0.74$~\si{Mm^3}\ shortfall is consistent with the DEM's own small
residual negative bias above. Across the two other reservoirs carrying an independent
volume reference at comparable A/P, Rosamarina and Poma, the reconstruction's band
capacity change differs from that reference by a mean absolute $10.1\%$, which is the
uncertainty to budget for at Garcia's A/P and above. We report the raw, uncorrected
estimate throughout rather than propagating a correction fitted at a single reservoir.

\begin{figure}[htbp]\centering
  \includegraphics[width=\linewidth]{garcia_validation_deep.png}
  \caption{Garcia validation against the December-2025 echo-sounder survey, using the
  curated official quota-area-volume table derived from it, the same reference as
  Tables~\ref{tab:hyps} and~\ref{tab:capacity}. (left) Volume--elevation curves for the
  design curve, the survey and the SAR DEM. (right) Storage loss against design at four
  levels, with the share of the surveyed loss the SAR DEM recovers.}
  \label{fig:garcia}
\end{figure}

\subsubsection{Validation against recent bathymetric curves}
\label{sec:res_curves}
Six reservoirs have a recent official area--volume curve as a secondary reference
(Table~\ref{tab:capacity}, rightmost column; Fig.~\ref{fig:aevgrid}). The comparison is
informative less for the raw size of any single gap than for its pattern across A/P. At high
A/P, SAR and the official curve agree closely. Poma's official curve shows essentially no net
loss ($-0.5\%$ band, $+2.9\%$ total), while the SAR DEM shows a modest $-12.8\%$ band deficit
attributable to the reconstruction's own low-elevation bias (Section~\ref{sec:res_garcia})
rather than real sedimentation, a useful negative-control case. Rosamarina and Arancio show
a similar pattern at a smaller offset (SAR $-15.1\%$/$-15.5\%$ versus survey/curve
$-7.3\%$/$-3.7\%$ band). At low A/P the gap widens sharply: Castello ($-30.0\%$ vs.\
$-23.5\%$), Olivo ($-45.3\%$ vs.\ $-3.6\%$) and Nicoletti ($-40.2\%$ vs.\ $-6.1\%$) all show
SAR over-estimating the true band loss by a wide margin, Olivo and Nicoletti most severely.
This is the same A/P-ordered pattern already visible in Section~\ref{sec:res_change}:
reconstruction reliability, not just apparent capacity change, degrades as A/P falls.

\subsubsection{Cross-sensor validation: PlanetScope optical reconstruction}
\label{sec:res_xsensor}
Reconstructing the bathymetry from PlanetScope 3~m optical masks (Section~\ref{sec:xsensor};
7--58 masks per reservoir) and comparing to the Sentinel-1 DEM, the two independent sensors
agree closely across the observable band at medium/high A/P: the per-pixel elevation RMSE is
$1.25$~m (Pozzillo, A/P~240), $1.65$~m (Poma, A/P~190) and $1.72$~m (Rosamarina, A/P~187),
of the same order as the Garcia \emph{field} RMSE ($2.67$~m) (Fig.~\ref{fig:planet}). At
Pozzillo, the optical masks are paired with SWOT rather than the gauge
(Section~\ref{sec:res_inputvalid}), for consistency with its primary reconstruction and, with
both series offset-corrected and cleaned of the same confirmed bad dates
(Section~\ref{sec:res_inputvalid}), a further, independent corroboration on top of the
SAR-area evidence that already motivated the SWOT-over-gauge choice there. Expressed as total
band volume, the optical and SAR reconstructions differ by $2.7$~\si{Mm^3} at
\emph{Pozzillo} ($23.5$ vs $20.8$) and $2.5$~\si{Mm^3} at \emph{Rosamarina} ($35.4$ vs $38.0$)
These gaps are consistent with, not larger than, the per-pixel RMSE integrated over each
reservoir's observable area, and they give these two reservoirs, which lack a modern field
survey, an independent near-truth on the same footing as Garcia's echo-sounder.

Rosamarina's and Pozzillo's comparison bands are each narrower than their full
SAR-reconstructed range, but no longer for the same reason. New PlanetScope acquisitions from
the second half of 2025 into 2026 (Section~\ref{sec:data}) close what was previously an
archive gap at the \emph{top} of Rosamarina's range: the optical band now reaches
$162.4$~m, matching the SAR reconstruction's own $162.3$~m ceiling; a modest gap remains
instead at the \emph{bottom} ($145.1$ vs.\ the SAR floor at $140.6$~m), where even the
earliest available PlanetScope acquisitions do not reach quite as low as the deepest
SAR-observed drawdown. Pozzillo received no new acquisitions in this round: its optical masks
(once paired to SWOT) still reach only $350.3$~m, a real $\sim$$4.6$~m gap against the SAR
reconstruction's own $354.9$~m ceiling. Rosamarina also illustrates why the water-level source
matters within its shared window: its in-situ gauge was stuck for $\sim$19 months across the
2024--2026 window (Section~\ref{sec:res_inputvalid}); substituting SWOT, which like the
optical masks records the real drawdown, restores a physical bathymetry and the reported
cross-sensor agreement. The two reconstructions diverge sharply (RMSE $8.05$~m) only at
\emph{Ancipa} (A/P~90); tellingly, this divergence \emph{grew}, rather than shrank, once two
new PlanetScope acquisitions (2024-11-14, 2024-11-25) were added specifically to close a
gap in its archive that fell inside the SAR reconstruction's own drawdown window
(2024-05-24 to 2024-12-19). The RMSE was $6.06$~m with 5 masks, before these two dates were
available. Extending the optical record deeper into Ancipa's low-A/P, narrow basin did not
resolve the mismatch; it sharpened it, indicating the divergence is a genuine, geometry-driven
reconstruction limit at low A/P rather than an artifact of a temporal sampling gap that more
data would fix. Because optical NDWI and SAR backscatter fail under
orthogonal conditions, agreement between them elsewhere is strong evidence the reconstruction
is not dominated by sensor-specific error.

\begin{figure}[htbp]\centering
  \includegraphics[width=\linewidth]{planet_vs_sar.png}
  \caption{Cross-sensor check: storage curves (area against volume) from the PlanetScope
  (optical, 3~m) and Sentinel-1 (SAR) bathymetric reconstructions, over the common
  observable band, for the four reservoirs with PlanetScope coverage.}
  \label{fig:planet}
\end{figure}

Agreement between two reconstructions is not the same as either being right, so the two
are also compared against the original design curve in the space the reconstruction
actually produces, volume against elevation (Fig.~\ref{fig:planetvol}). Each curve is
drawn only over the band its own sensor observed, which makes the differing ceilings
part of the result rather than a hidden assumption. At \emph{Pozzillo} the SAR
reconstruction tracks the design curve closely to its own ceiling at $354.9$~m and
recovers $46.24$~\si{Mm^3}, while the optical reconstruction stops at $350.3$~m with
$27.07$~\si{Mm^3}, so the difference between them is one of coverage rather than of
disagreement where both see. At \emph{Ancipa} the ordering reverses. The optical
reconstruction reaches $946.5$~m and $23.05$~\si{Mm^3}, essentially on the design curve
($22.21$~\si{Mm^3} at the same level), whereas the SAR reconstruction recovers only
$13.42$~\si{Mm^3} over a nearly identical band. Ancipa is the lowest-A/P reservoir
carrying an optical check, and this is the same under-reconstruction its capacity
numbers show in Section~\ref{sec:res_curves}, now visible against a second sensor rather
than inferred from a curve. At \emph{Poma} and \emph{Rosamarina}, where an updated
official curve exists, the two reconstructions bracket it and stay within a few
\si{Mm^3} of each other over the shared band.

\begin{figure}[htbp]\centering
  \includegraphics[width=\linewidth]{planet_volume_elevation.png}
  \caption{Volume--elevation curves from the PlanetScope and Sentinel-1 reconstructions
  against the design curve, and the updated official curve where one exists, for the four
  reservoirs with PlanetScope coverage. Volumes are relative to a common floor per
  reservoir, and each reconstruction is drawn only over the band it observes, so the
  extent of each curve is itself the observability limit of that sensor.}
  \label{fig:planetvol}
\end{figure}
